\section[Decisions]{Architecture decision records}
\sectionframe{05}{Architecture decision records}
  {Why does the project work this way, and what if we change our minds?}
  {\faBalanceScale}

% ------------------------------------------------------------ what / when
\begin{frame}[source={Nygard (2011), Documenting architecture decisions · Pocock (2026), mattpocock/skills v1.3.1, domain-modeling}]{What an ADR is, and when to write one}
  \begin{tikzpicture}[x=1mm,y=1mm,
      crit/.style={inner sep=0pt,font=\footnotesize,text=udsdark,align=center},
      item/.style={anchor=base west,inner sep=0pt,font=\footnotesize,text=udsdark}]
    % the record itself: Nygard's format
    \node[anchor=north west,inner sep=0pt,font=\footnotesize,text=udsdark] at (0,0)
      {\textbf{One decision per record:}\enspace title · context · decision · status ·
       consequences; numbered; superseded records stay};
    % write one only where all three hold; each label sits in its circle's own part
    \coordinate (A) at (17,-22); \coordinate (B) at (35,-22); \coordinate (C) at (26,-37.6);
    \foreach \p in {A,B,C}
      \filldraw[draw=udsblue,line width=0.7pt,fill=udsblue,fill opacity=0.09] (\p) circle (15);
    \node[crit] at (11.5,-19.5) {hard to\\reverse};
    \node[crit] at (40.5,-19.5) {surprising\\without\\context};
    \node[crit] at (26,-44) {a real\\trade-off};
    \node[inner sep=0.4mm,font=\small\bfseries,text=udsred] (adr) at (26,-27.2) {ADR};
    % where all three usually hold in research code, level with the circles' top
    \node[anchor=north west,inner sep=0pt,font=\footnotesize\bfseries,text=udsdark] at (70,-7)
      {typical in research code};
    \foreach \t [count=\i] in {the split unit and the evaluation protocol,merged or remapped classes,
        {units, scales, normalisation},the metric that decides what is ``better''}
      \node[item] at (72,-9.2-\i*5.4) {\textcolor{udsred}{\faAngleRight}\enspace\t};
    \draw[udsred,line width=0.6pt] (adr.east) to[out=0,in=180] (69.4,-21.7);
    \node[anchor=base west,inner sep=0pt,font=\scriptsize,text=udsgrey] at (72,-41)
      {no ADR: plot styling, renames};
  \end{tikzpicture}
  \note{\textbf{20:30--21:30}\par
    An architecture decision record keeps one decision with its reasons.\par
    Nygard's format: title, context, decision, status and consequences, including the bad
      ones.\par
    Records are numbered in sequence, numbers are never reused, and a superseded record stays in
      the repository.\par
    A useful test, from the skills of section 6: write one only when the decision is hard to
      reverse, would surprise a newcomer, and came from a real trade-off.\par
    In research code that is usually about what a number means: the split unit, class
      definitions, units and normalisation, the metric.\par
    Plot styling and renames do not need one.}
\end{frame}

% ------------------------------------------------------------- ADR example
\begin{frame}[fragile]{ADR-0003: specimen-level splits}
  \begin{columns}[T,onlytextwidth]
    \begin{column}{0.58\textwidth}
\begin{codefile}[basicstyle=\ttfamily\fontsize{7}{8.4}\selectfont]{docs/adr/0003-specimen-level-splits.md}
# ADR-0003: Specimen-level splits
Status: accepted. Supersedes ADR-0001 (random patch split).
## Context
Micrographs of one specimen share preparation and etching;
several show the same area. A random patch split puts this
material in training and validation.
## Decision
Specimens are assigned to validation folds; each patch
inherits its specimen's fold. Folds are saved with seed
and manifest hash and checked for overlap.
## Considered options
Random patch split; split by micrograph. Both leak.
## Consequences
Few independent units: report every fold and its spread.
One specimen per state: a held-out specimen is a new state.
\end{codefile}
    \end{column}
    \begin{column}{0.39\textwidth}
      \begin{tikzpicture}[x=1mm,y=1mm]
        \node[artifact,anchor=north west,text width=50mm,align=left,font=\scriptsize] (a1) at (0,0)
          {\texttt{0001-random-patch-split.md}\\[0.4mm]
           \textcolor{udsred}{Status: superseded by ADR-0003}};
        \node[artifact,anchor=north west,text width=50mm,align=left,font=\scriptsize,draw=udsblue,fill=bluetint] (a3) at (0,-16)
          {\texttt{0003-specimen-level-splits.md}\\[0.4mm]
           \textcolor{udsblue}{Status: accepted}};
        \draw[flow] (a3.north) -- node[lbl,right]{supersedes} (a1.south);
        \node[anchor=north west,text width=54mm,font=\scriptsize,inner sep=0pt] at (0,-27.5)
          {ADR-0001 stays: it shows what was tried and why it changed.};
      \end{tikzpicture}\par\vskip3mm
      \begin{panel}[fontupper=\footnotesize,colback=bluetint,borderline west={1.2mm}{0pt}{udsblue},left=3.5mm]
        An agent chasing a better score finds ADR-0003 through \fn{AGENTS.md} before it brings back a
        random split.
      \end{panel}
    \end{column}
  \end{columns}
  \note{\textbf{21:30--22:30}\par
    The record for our running example.\par
    The context explains why a random patch split misleads here; the decision is short and
      testable; the considered options show that splitting by micrograph was rejected for the
      same reason.\par
    The consequences are honest: few specimens mean few independent units, and with one specimen
      per material state every held-out specimen is also a new state.\par
    ADR-0001 is not deleted; its status points to its successor.\par
    An agent tuning for a better score finds this record before it undoes the decision.}
\end{frame}

% ------------------------------------------------------------- file roles
\begin{frame}[source={OpenAI (2026), Custom instructions with AGENTS.md · Pocock (2026), mattpocock/skills v1.3.1: code-review; v1.3 changelog}]{Six files, six jobs}
  % the six files of sections 1, 4 and 5 placed on an agent session, at the moment each one is read
  \begin{tikzpicture}[x=1mm,y=1mm,
      chip/.style={draw=udsgreylight,fill=white,rounded corners=0.6mm,inner xsep=1mm,inner ysep=0.8mm,
        anchor=north west,font=\ttfamily\fontsize{7.4}{8.4}\selectfont,text=udsdark,line width=0.5pt,
        path picture={\draw[udsblue,line width=0.9pt] (path picture bounding box.north west)
          -- (path picture bounding box.north east);}},
      always/.style={chip,draw=udsblue,line width=0.9pt},
      ondemand/.style={chip,dashed,draw=udsdark,path picture={}},
      reviewstep/.style={chip,draw=udsblue,fill=bluetint,path picture={},font=\footnotesize},
      say/.style={anchor=north west,inner sep=0pt,font=\scriptsize,text=udsdark,align=left,text width=30mm,
        execute at begin node={\hyphenpenalty=10000\raggedright}},
      drop/.style={udsgreymid,line width=0.5pt},
      dot/.style={circle,fill=udsdark,inner sep=0pt,minimum size=1.8mm}]
    % AGENTS.md points to the glossary and the ADRs
    \node[always] (ag) at (0,-4.6) {AGENTS.md};
    \node[ondemand] (gl) at (44,-4.6) {GLOSSARY.md};
    \node[ondemand] (ad) at (78,-4.6) {docs/adr/};
    \draw[drop,dashed] (ag.north) |- (ad.north |- 0,-1.6);
    \foreach \n in {gl,ad} \draw[drop,dashed,-{Stealth[length=1.6mm,width=1.2mm]}] (\n.north |- 0,-1.6) -- (\n.north);
    \node[anchor=south,inner sep=0.5mm,font=\scriptsize,text=udsgrey,fill=white] at (27,-1.6) {points to};
    \node[say] at ([yshift=-1.2mm]ag.south west) {\textcolor{udsblue}{\bfseries every session}\\ commands, boundaries, pointers};
    \node[say] at ([yshift=-1.2mm]gl.south west) {\textbf{when pointed to}\\ domain terms};
    \node[say] at ([yshift=-1.2mm]ad.south west) {\textbf{when pointed to}\\ one decision and its reasons};
    % read in review
    \node[chip] (cs) at (112,-4.6) {CODING\_STANDARDS.md};
    \node[say] at ([yshift=-1.2mm]cs.south west) {\textbf{in review}\\ how code is written here};
    % the session, left to right, and where each file enters it
    \draw[flow,line width=0.9pt] (0,-27) -- (143.4,-27);% the tip stays inside the text width
    \node[anchor=south west,inner sep=0.6mm,font=\scriptsize,text=udsgrey] at (12,-27) {an agent session};
    \foreach \n in {ag,gl,ad,cs} {\node[dot] at (\n.south |- 0,-27) {}; \draw[drop] (\n.south |- 0,-21.6) -- (\n.south |- 0,-27);}
    % the task's own specification comes with it; the review checks the code against the standards
    \node[chip,font=\footnotesize] (sp) at (0,-32) {task specification};
    \node[say] at ([yshift=-1.2mm]sp.south west) {\textbf{with the task}\\ objective, acceptance criteria};
    \draw[drop] (ag.south |- 0,-27) -- (ag.south |- sp.north);
    \node[reviewstep] (cr) at (112,-32) {\faSearch\enspace review};
    \node[say] at ([yshift=-1.2mm]cr.south west) {a person or an agent checks the code against it};
    \draw[drop] (cs.south |- 0,-27) -- (cs.south |- cr.north);
    % read like any other file
    \node[chip,draw=udsgreylight,text=udsgrey,path picture={}] (rd) at (44,-34) {README.md};
    \node[say,text=udsgrey,text width=40mm,anchor=west] at ([xshift=2mm]rd.east) {for people; read like any other file};
  \end{tikzpicture}
  \note{\textbf{22:30--23:30}\par
    All six files on one agent session, at the moment each one is read.\par
    AGENTS.md at the start of every session; the glossary and the ADRs when AGENTS.md points to
      them.\par
    The task specification of section 1 comes with the task; the coding standards are read when
      the change is reviewed, by a person or an agent.\par
    The README is for people; the agent reads it like any other file.\par
    Keep each fact in one place and link to it: copies go stale.}
\end{frame}

% ------------------------------------------------ the request, revisited
\begin{frame}{The same request, in a prepared repository}
  \begin{tikzpicture}[x=1mm,y=1mm]
    % the cards fill the text width exactly; strokes on their edges must not widen the picture
    \useasboundingbox (0,0) rectangle (144,-54);
    \specterminal{0}{the request}
    \node[spec request] at (term) {> Add validation folds by specimen to the training CLI, as decided in ADR-0003.};
    \node[anchor=west,inner sep=0pt,font=\scriptsize\bfseries,text=udsdark] at (0,-15.8)
      {where the rest comes from};
    % the same five cards as on slides 4 and 5: the objective is the request itself, every other
    % part now comes from files
    \speccard{0}{-19}{-44.4}{Objective}{darktint}
    \foreach \k/\l in {1/Context,2/Constraints,3/Acceptance,4/Verification}
      {\specfilecard{\k}{-19}{-44.4}{\l}{darktint}}
    \node[spec chip,draw=termbg,fill=termbg,text=termfg,path picture={}] at (b0) {\faTerminal\ the request};
    % context
    \node[spec chip] (c1) at (b1) {GLOSSARY.md};
    \node[spec note] (c1s) at ([yshift=-0.8mm]c1.south west) {specimen, fold};
    \node[spec chip] (c2) at ([yshift=-2.2mm]c1s.south west) {ADR-0003};
    \node[spec note] at ([yshift=-0.8mm]c2.south west) {why; how folds are saved};
    % constraints
    \node[spec chip] (k1) at (b2) {AGENTS.md};
    \node[spec note] at ([yshift=-0.8mm]k1.south west) {boundaries};
    % acceptance
    \node[spec chip] (a1) at (b3) {ADR-0003};
    \node[spec note] (a1s) at ([yshift=-0.8mm]a1.south west) {one fold per specimen};
    \node[spec chip] (a2) at ([yshift=-2.2mm]a1s.south west) {CODING\_STANDARDS.md};
    \node[spec note] at ([yshift=-0.8mm]a2.south west) {seed, input checks};
    % verification
    \node[spec chip] (v1) at (b4) {AGENTS.md};
    \node[spec note] at ([yshift=-0.8mm]v1.south west) {fast check, report};
    % what the files leave open, the agent has to ask: the clarifying step of the next section
    \node[anchor=west,inner sep=0pt,font=\scriptsize,text=udsred] at (0,-49.4)
      {\faQuestionCircle\enspace still open: what no file says, e.g.\ the default number of folds};
  \end{tikzpicture}
  \note{\textbf{23:30--24:50}\par
    Back to the request from the first section.\par
    With these files in place, one sentence in the project's language is enough.\par
    Four of the five parts now come from files; only the objective is in the request.\par
    The glossary says what a specimen and a fold are.\par
    ADR-0003 says why folds go by specimen and how they are saved.\par
    AGENTS.md carries the boundaries, the test command and what to report.\par
    CODING\_STANDARDS.md makes seeds and input checks part of every review.\par
    What the files leave open, here the default number of folds, is what the clarification step
      in the next section asks.\par
    This works only while AGENTS.md points to the files and they are kept current.}
\end{frame}
